\documentclass[12pt]{article}
\usepackage[utf8]{inputenc}
\usepackage[T1]{fontenc}
\usepackage{hyperref}
\title{Northern Kankanay Language Resources for K-12 Curriculum Development}
\author{Compiled by Claude Sonnet 5.5 (claude-sonnet-5-5) for Andr\'as Kornai}
\date{}
\begin{document}
\maketitle

\section{Introduction}
Northern Kankanay (ISO 639-3 xnn; Glottocode nort2877) is an Austronesian language of the Philippines, spoken in western Mountain Province, southeastern Ilocos Sur and southern Abra. Other names include Sagada Igorot and Western Bontoc. A Webonary page cites 94,100 speakers (2005 UNSD); the brief gives about 70,000. The search found a dictionary and a grammar reference, but no school materials.

\section{Primary Educational Resources}
No primers, textbooks or school programmes were found in search results.

\section{Linguistic and Reference Materials}
\begin{itemize}
\item Northern Kankanay--English dictionary (SIL, compiled by Judy Wallace, 2018, still under development), with an overview that mentions a completed grammar; Latin script: \url{https://www.webonary.org/northernkankanay/overview/introduction/}
\item Kankanaey language (Wikipedia), which treats Northern Kankanaey as separate from Kankanaey: \url{https://en.wikipedia.org/wiki/Kankanaey_language}
\item Separate Kankanaey--English dictionary on Webonary, documenting the Poblacion, Kibungan, Benguet dialect. This is a different variety and should not be mixed with xnn: \url{https://www.webonary.org/kankanaey/?p=159999}
\end{itemize}
I did not find the title or publisher of the Northern Kankanay grammar.

\section{Religious and Community Resources}
None found in search results.

\section{Digital and Audio Resources}
The Webonary dictionary is the only digital resource found. I did not verify whether it includes audio.

\section{Curriculum Development Context}
A healthy speaker base and a written dictionary give a starting point for mother-tongue materials. Developers should confirm with Philippine education authorities and local communities which materials exist, since none turned up here. They should also check carefully that Kibungan-dialect material is not treated as Northern Kankanay.
\end{document}